\subsection{Support at the fixed level places}
\label{subsec:fixed-level-support}

At the fixed level places the scalar in the unfolded integral is determined by the support of the two translated test functions. At a fixed level place the first flat map
has support in the product of the Jacobi subgroup and the
specified level. After contraction, the second flat map
has support in the product of the Bessel subgroup and its
specified level. These statements concern the translated
test functions on the chosen level lattices. We prove them
for the split model of the fixed level. In the formulas below
the conductor of the additive character \(\psi\) is \(O\),
the integer ring, and its additive Haar measure \(\nu\)
gives \(O\) mass one.

\begin{lemma}[Support of the split position functional]\label{lem:split-position-support}
At a fixed split level place with the symmetric chain and
its standard iterated type function \(\Phi\), let \(c\) be the number of numbered target coordinates. For \(1\le p\le c\), let \(P_p\) integrate the test function against the prescribed additive character over the free entries of the numbered slice at \(p\). With \(R_p\) the position subgroup and \(J\) the prescribed level, the support satisfies
\[
 \{h:P_p(h\Phi)\ne0\}\subseteq R_pJ.
\]

\end{lemma}

\begin{proof}
\proofstep{Record the integral flag coordinates and their residue support.}
Write \(k\) for the local field, \(\varpi\) for a uniformizer,
and \(\Phi\) for the standard iterated type function at a stage
of a tuple. There are \(c\) numbered target flag coordinates
and \(c+1\) source flag coordinates, with ordinary coordinates
in the remaining block. With the ordinary block first, the target and source frame
matrices have the form
\[
 D=I_{\rm ord}\oplus\operatorname{diag}(\varpi,\ldots,\varpi^c),
 \qquad D'=I_{\rm ord}'\oplus
             \operatorname{diag}(\varpi,\ldots,\varpi^{c+1}).
\] The field coordinates of
the test matrix \(X\) are read as \(D^{-1}XD'\). In these
coordinates the level \(J\) is integral. The entries on
and below the prescribed parts of the flag diagonal agree
with the identity modulo \(\varpi\). The iterated test keeps
the retained columns of a matrix in that residue pattern;
its other coordinates are Schwartz coordinates. Let \(M\) be the corresponding rectangular matrix space. In the
original coordinates the residue condition reads
\[
 \operatorname{supp}\Phi\subseteq D M(O)D'^{-1},\qquad
 \Phi(X+DZD'^{-1})=\Phi(X)\quad(Z\in\varpi M(O)).
\]

At a numbered position \(1\le p\le c\), the operator \(P_p\)
integrates \(\Phi(X)\psi(-\sum_i X_{i,i+1})\) over the free
entries of \(X\). In the character sum the two indices run
in the numbered flag blocks. The ordinary block is the
argument of \(P_p\Phi\), and the numbered diagonal entries
of the affine slice are one.

Write \(R_p\) for the group
generated by the ordinary general linear block and the
target unipotent for this slice. The latter has scalar
covariance under \(P_p\), and the ordinary block acts by
invertible changes of the argument and a nonzero scalar.
Let \(m(h)\) be the product of the determinant character
and modulus for a target translation. Its Weil action satisfies
\[
 (h\Phi)(X)=m(h)\Phi(h^{-1}X),\qquad m(h)\ne0.
\]
In particular the scalar does not alter the translated support. These two covariance
statements follow by changing the free variables. They
apply to the same numbered slice before and after translation.
We claim that
\begin{equation}\label{eq:position-level-support}
                   P_p(h\Phi)\ne0\quad\Longrightarrow\quad h\in R_pJ .
\end{equation}

For the first reduction use Iwasawa on \(hD\), with weights
\(2p-1\) on the ordinary block and \(2i\) on numbered
line \(i\). This writes \(h=n M w D^{-1}\), where \(w\) is
integral with integral inverse, \(n\) is in the target
unipotent of \(R_p\), and \(A=M^{-1}\) has block form
\[
          A=\begin{pmatrix}P&Q\\0&S\end{pmatrix},
       \qquad S=\operatorname{diag}(s_i),\qquad Q_i=0\ \ (i\ne p).
\]
Here \(Q_i\) is column \(i\) of \(Q\). This form follows by
taking the Levi for the order of the numbered lines with
the ordinary block tied to line \(p\), and using the finer
Iwasawa order inside that grade.

Removing \(n\) preserves
nonvanishing. The integral change \(w\in\mathrm{GL}(O)\) preserves both
\(M(O)\) and \(\varpi M(O)\). Thus the transformed coordinates
give the necessary support and the translation lattice
\[
 AXD'\in M(O),\qquad
 Z\in A^{-1}\varpi M(O)(D')^{-1}.
\] At a nonzero reading the
numbered diagonal entries give \(s_i\varpi^i\in O\); column
\(p\) has zero ordinary part and gives \(Q_p\varpi^p\in O\).

For the inverse diagonal condition, let \(y\) range over
\(s_i^{-1}\varpi^{-i}O\). Translate the integration variables
by \(Z=A^{-1}E\), with \(E\) supported at numbered position
\((i,i+1)\) with entry \(ys_i\). Then \(AZD'\) is divisible
by \(\varpi\). The translated variables are free, including
the possible ordinary part of column \(p+1\) when \(i=p\).
The character is multiplied by \(\psi(-y)\).

If the reading is \(I\ne0\), translation invariance gives
\[
 I=\psi(-y)I\quad(y\in s_i^{-1}\varpi^{-i}O)
 \quad\Longrightarrow\quad s_i^{-1}\varpi^{-i}O\subseteq O .
\]
Here \(O\) is the conductor of \(\psi\). Together with the numbered
diagonal condition this proves \(s_i\varpi^i\in O^\times\). With
\(\gamma=P^{-1}\) in the ordinary block,
\(A\operatorname{diag}(\gamma,1)D\) is integral with
integral inverse: its diagonal blocks are the identity and
a diagonal unit matrix, and its upper block is
\(Q\operatorname{diag}(\varpi^i)\).
This gives the factorization required for the first reduction:
\[
 h=r u,\qquad r\in R_p,\qquad u\in D\mathrm{GL}(O)D^{-1}.
\]

For the second reduction it suffices to treat this \(u\), using
the covariance of \(P_p\) under \(R_p\). Nonvanishing gives an
affine-slice value of the translated test. In frame
coordinates this gives a \(\kappa\) in the prescribed residual
level such that \(g=D^{-1}uD\,\kappa\) agrees with the scaled
affine matrix on the retained columns modulo \(\varpi\).
Write \(G=D^{-1}uD\) and let \(Z\) consist of the retained
columns of the scaled affine matrix. The nonzero reading of the
previous type says
\[
             G^{-1}Z\equiv\binom{g_0}{0}\pmod{\varpi},
\]
where \(g_0\) belongs to the previous level. In the reordered
frame \(\kappa=\operatorname{diag}(g_0,1)\) is in the next level.
Thus the retained columns of \(G\kappa\) agree with \(Z\) modulo
\(\varpi\).
The affine zeros imply the following two
conditions. The lower entries in the retained initial
ordinary columns, in rows \(i\ge p\), vanish modulo
\(\varpi\).

In every numbered column the entries at or below
that column's grade agree with the identity modulo
\(\varpi\). Replace those entries by their exact zero and
one values, fixing all other entries. Writing \(\widetilde g\) for the changed matrix, one has
\(g^{-1}\widetilde g\equiv1\pmod\varpi\), so this change is made
inside the prescribed right level. Its lower numbered square
block \(S_0\) is upper triangular unipotent. If \(E_0\) is
the lower rectangular block, then
\[
     (E_0\ S_0)
         \begin{pmatrix}1&0\\-S_0^{-1}E_0&1\end{pmatrix}
          =(0\ S_0).
\]
The multiplying block is zero on the retained initial ordinary
columns because \(E_0\) is zero there. On the other columns this
shear belongs to \(J\) by the iterated level conditions.
The result belongs to the frame reading of \(R_p\) up to
principal congruence. Each right factor used belongs to
the prescribed level. This proves \eqref{eq:position-level-support}
with that level rather than the full lattice stabilizer.
\label{cor:period-level-support-1}

\paragraph{Exchange of adjacent positions.}
The position calculation uses the symmetric chain of even flag length, so the midpoint and full-length parameters satisfy
\[
 s=2\rho,\qquad 1\le i\le\rho\quad\text{in the first half}.
\] For a move from position \(i\) to \(i+1\) with
\(i\le\rho\) and \(i+1\le c\), the two affine integrals have the
same free variables except for an ordinary row and an ordinary
column. The integral for \(i+1\) frees the entries in numbered
target row \(i\) and the ordinary source columns. The integral
for \(i\) instead frees the entries in numbered source
column \(i+1\) and the ordinary target rows. The compensating blocks lie in the kernel of the additive character in the two affine integrals: if \(b\) denotes the change in its argument in either block,
\[
 \psi(a+b)=\psi(a)\psi(b)=\psi(a).
\]

The first block contributes
the indicator of \(\varpi^iO\) in each coordinate: the frame
divides these entries by \(\varpi^i\), and the permitted
integral row addition preserves the test. This addition has
unipotent determinant one and leaves the residual type
character trivial on the added row. The second contributes
the indicator of \(\varpi^{-i}O\): the frame multiplies the
column by \(\varpi^{i+1}\), and these entries are prescribed
zero modulo \(\varpi\). In each free entry the two tests reduce respectively to
\[
 \varpi^{-i}x\in O\ \Longleftrightarrow\ x\in\varpi^i O,
 \qquad \varpi^{i+1}y\in\varpi O\ \Longleftrightarrow\
                       y\in\varpi^{-i}O .
\]
After removing these indicators both restrictions give the same
test with the free block zero.

For the row this uses the iterated row pattern, for which
position \(i\le\rho\) permits the shear after the initial
row has been set to zero. For the column it uses the
periodicity modulo \(\varpi\) of the frame reading.
Thus the ratio of the two operators is the positive scalar
\[
     \frac{\nu(\varpi^iO)^{d_{\mathrm{src}}}}
          {\nu(\varpi^{-i}O)^{d_{\mathrm{tgt}}}},
\]
where the exponents are the respective ordinary block
dimensions. Iteration compares the operator at any
position needed here to the operator at position one.
Empty ordinary blocks contribute factors one. The ratio
depends on the place, dimensions and position and is common
to the tuples, which use the same lattice exponents and
the same measure.
\end{proof}

\paragraph{The two flat maps.}
Use Lemma~\ref{lem:split-position-support} at each position of
the symmetric chain.
In the Bessel to Fourier--Jacobi step the flat column map
at position \(p\) with extra column \(y\) reads
\(P_p(E_y^{-1}h\Phi)\), with \(E_y\) the upper shear that
adds the extra column into the ordinary block and has
determinant one. This follows by translating the affine
matrix by the added column before taking the Fourier
integral. Equation~\eqref{eq:position-level-support} yields
\[
  P_p(E_y^{-1}h\Phi)\ne0
     \ \Longrightarrow\ E_y^{-1}h\in R_pJ
     \ \Longrightarrow\ h\in E_yR_pJ .
\]
In the split frame the product \(E_y r\) lifts to the
Jacobi group. The right factor lifts to precisely the
target level at that stage because the other split block
is inverse adjoint. This gives the first support statement.

For the value calculation at the identity, the middle
column and row give pure factors. With \(c=2a\) in a double
step, the first flat column in its middle variable is
a position-one entry times a positive constant times
\(\mathbf1_{\varpi^{-a-1}O^d}\). The following flat row is
its own position-one entry times a positive constant times
\(\mathbf1_{\varpi^{a+1}O^d}\), with \(d\) the middle
ordinary dimension. For the added entries \(y\) in numbered column \(a+2\), the
source frame and its zero residue condition give
\[
  \varpi^{a+2}y\equiv0\pmod{\varpi O^d}
         \quad\Longleftrightarrow\quad y\in\varpi^{-a-1}O^d .
\]

In this
range the addition is a translation by the period lattice
of the test function.
For the row the ordinary-to-numbered conversion is the
source shear with the direct range \(\varpi^{a+1}O^d\).
It has determinant one and belongs to the source level
with trivial character at these row coordinates. Outside
the range, the affine row test is zero by the dual lattice
condition. The remaining operator in each case is put at
position one by the exchange calculation. The position-one
normalization restores the positive Haar volume of its
radical stabilizer.

To compare the two position-one entries, let \(C_0,R_0\) denote those entries and \(c_C,c_R>0\) their pure-map factors. The common lattice is
\[
 L=\varpi^{-a-1}O^d .
\] The column and row functions are
\(c_C C_0\mathbf1_L\) and \(c_R R_0\mathbf1_{L^\perp}\).
For the additive character \(\psi\) of the middle model,
\[
 \widehat{\mathbf1_L}(z)=\int_L\psi(\langle x,z\rangle)\,dx
             =\nu^d(L)\mathbf1_{L^\perp}(z),
 \qquad L^\perp=\varpi^{a+1}O^d.
\]
Indeed changing \(x\) by \(x_0\in L\) multiplies the integral
by \(\psi(\langle x_0,z\rangle)\); it is therefore zero unless
this equals one for every \(x_0\). At the identity the
contraction is \(c_Cc_R C_0R_0\nu^d(L)\nu^d(L^\perp)\).

For the translated second test let \(R_h(z)\) be its row
function. The contraction is
\[
 c_C C_0\nu^d(L)\int_{L^\perp}R_h(z)\,dz .
\]
On \(L^\perp\) the source shear belongs to the level and
commutes with the left target translation; thus \(R_h(z)\)
is the position reading of \(h\Phi\) multiplied by its
nonzero shear character and pure-map factor. A nonzero value
forces that reading to be nonzero, whence \(h\in R_pJ\) by
\eqref{eq:position-level-support}. Here \(R_p\) lifts to the
target Bessel subgroup. Only the transformed column confines
this integral. The support of the translated second test can
be larger.

Both deductions use the split Weil
model and the compatible Jacobi and Bessel frames of
the same step. In the second frame let \(e_\xi\) denote the image of the
negative line generator under the fixed line identification. The
normalization in this split model is
\[
 \langle e_\xi,e_\xi\rangle=-\xi^2 .
\]
\label{cor:period-level-support-2}

\paragraph{Use at all the fixed places.}
Let \(S\) denote the fixed places and write the translation
as \(h=\prod_{v\in S} h_v\). The operators on distinct local
Weil factors commute. For each \(v\), let \(T_v\) be the local contraction or position map on the corresponding Weil factor. Applied to a finite expansion \(\sum_a c_a\bigotimes_v\phi_{a,v}\), the tensor product of these maps gives
\[
  \sum_a c_a\bigotimes_v T_v(\phi_{a,v}).
\]

If one \(T_v\) vanishes on the full local test span, each summand
vanishes, independently of the other places.
Thus a nonzero global value implies membership in the local cell at every \(v\), giving
\(h\in\prod_{v\in S} R_vJ_v\) in the corresponding local
notation. Multiplying the local Jacobi or Bessel factors
and the local levels gives the fixed-level cell, using the
specified product level. At the identity the common
multiplier is the product of the positive local constants.
The terminal two steps use the rank-one line identification
and compact terminal pairing for this common-value assertion;
the double-step support computation concerns the
nonterminal steps. These are the cell integration inputs.


\label{cor:period-level-support-3}
